 \bta{2013}
 \begin{enumerate}
 	%\renewcommand{\labelenumi}{\arabic{enumi}.}
 	% A(\Alph) a(\alph) I(\Roman) i(\roman) 1(\arabic)
 	%设定全局标号series=example	%引用全局变量resume=example
 	%[topsep=-0.3em,parsep=-0.3em,itemsep=-0.3em,partopsep=-0.3em]
 	%可使用leftmargin调整列表环境左边的空白长度 [leftmargin=0em]
 	\item
 	（共30分）
 	质量为 $\mu$ 的粒子在一个无限深球方势阱
 	$$
 	V(r)=\left\{\begin{array}{ll}
 		0, & r \leq a \\
 		\infty, & r>a
 	\end{array}\right.
 	$$
 	中运动。
 	\begin{enumerate}
 		%\renewcommand{\labelenumi}{\arabic{enumi}.}
 		% A(\Alph) a(\alph) I(\Roman) i(\roman) 1(\arabic)
 		%设定全局标号series=example	%引用全局变量resume=example
 		%[topsep=-0.3em,parsep=-0.3em,itemsep=-0.3em,partopsep=-0.3em]
 		%可使用leftmargin调整列表环境左边的空白长度 [leftmargin=0em]
 		\item
 		写出径向波函数 $R_{l}(r)$ 满足的方程（已知: $\nabla^{2}=\frac{1}{r} \frac{\partial^{2}}{\partial r^{2}} r-\frac{\hat{l}^{2}}{\hbar^{2} r^{2}}$ )。
 		\item 
 		求其中 $l=0$ 的归一化能量本征函数和能量本征值。
 		
 		
 		
 	\end{enumerate}
 	
 	
\banswer{
\begin{enumerate}
	%\renewcommand{\labelenumi}{\arabic{enumi}.}
	% A(\Alph) a(\alph) I(\Roman) i(\roman) 1(\arabic)
	%设定全局标号series=example	%引用全局变量resume=example
	%[topsep=-0.3em,parsep=-0.3em,itemsep=-0.3em,partopsep=-0.3em]
	%可使用leftmargin调整列表环境左边的空白长度 [leftmargin=0em]
	\item
	设波函数为$ \psi=R_{l}(r)Y_{lm}(\theta,\varphi) $，分离变量时产生的常数为$ l(l+1)\hbar^{2} $，得径向波函数为
	\[ 
	\frac{\mathrm{d}^{2}}{\mathrm{~d} r^{2}} R_{l}(r)+\frac{2}{r} \frac{\mathrm{d}}{\mathrm{d} r} R_{l}(r)+\left[\frac{2 \mu}{\hbar^{2}}(E-V(r))-\frac{l(l+1)}{r^{2}}\right] R_{l}(r)=0
	 \]
	可作变量替换$ R_{l}(r)=\chi_{l}(r) / r $，进一步化简为
	\[ 
	\chi^{\prime \prime}{ }_{l}(r)+\left[\frac{2 \mu}{\hbar^{2}}(E-V(r))-\frac{l(l+1)}{r^{2}}\right] \chi_{l}(r)=0
	 \]	
\item 
$l=0$ 时，$\chi''+k^2\chi=0$，其中 $k=\sqrt{2\mu E}/\hbar$。
原点不含 $\delta$ 势，物理波函数须在原点正则，因此 $\chi(0)=0$；
无穷高外壁要求 $\chi(a)=0$。这给出 $\chi=C\sin(kr)$、$ka=n\pi$。
由
\[
 1=\int_0^a|\chi|^2dr=|C|^2\frac a2
\]
得 $C=\sqrt{2/a}$。再乘 $Y_{00}=1/\sqrt{4\pi}$、除以 $r$，故在 $r>a$ 区域为零，
在 $0<r\le a$ 区域为
\[ 
\psi_{n}=\frac{1}{\sqrt{2a\pi} } \cdot \frac{\sin \frac{n\pi}{a}r}{r}
 \]	
能量本征值为
\[ 
E_{n}=\frac{n^{2}\pi^{2}\hbar^{2}}{2\mu a^{2}}
 \]
其中$ n=1,2,3,\cdots $
\end{enumerate}

	
}

 	
\item 
（共30分）考虑一质量为 $m$ 的自由粒子的一维运动。设初始 $(t=0)$ 时刻波函
数为 $\psi(x, 0)=\left(\frac{\alpha}{\pi}\right)^{1 / 4} e^{i k_{0} x-\alpha x^{2} / 2} $（$k_{0},\alpha$ 为实常数； $\int_{-\infty}^{\infty} d x e^{-a x^{2}}=\sqrt{\frac{\pi}{a}}, \quad a>0$）。
\begin{enumerate}
	%\renewcommand{\labelenumi}{\arabic{enumi}.}
	% A(\Alph) a(\alph) I(\Roman) i(\roman) 1(\arabic)
	%设定全局标号series=example	%引用全局变量resume=example
	%[topsep=-0.3em,parsep=-0.3em,itemsep=-0.3em,partopsep=-0.3em]
	%可使用leftmargin调整列表环境左边的空白长度 [leftmargin=0em]
	\item
	求 $t>0$ 时刻动量表象波函数 $\tilde{\Psi}(k, t)$ 及粒子动量几率分布 $\Pi(k, t)$ 。
	\item 
	求 $t>0$ 时刻波函数 $\Psi(x, t)$ 及粒子位置几率分布 $\mathrm{P}(x, t)$ 。
	\item 
	简述粒子动量几率分布 $\Pi(k, t)$ 及位置几率分布 $\mathrm{P}(x, t)$ 随时间演化的特性。
	
	
	
\end{enumerate}
\banswer{
以下用波数 $k=p/\hbar$ 为变量，约定 Fourier 核为 $e^{ikx}/\sqrt{2\pi}$。
归一化要求 $\alpha>0$。
\begin{enumerate}
\item 先对初态作 Fourier 变换，配方积分得
\begin{align*}
 \widetilde\Psi(k,0)
 &=\frac1{\sqrt{2\pi}}\left(\frac\alpha\pi\right)^{1/4}
   \int_{-\infty}^\infty e^{-\alpha x^2/2-i(k-k_0)x}dx\\
 &=(\pi\alpha)^{-1/4}e^{-(k-k_0)^2/(2\alpha)}.
\end{align*}
动量表象下 $H=\hbar^2k^2/(2m)$，所以
\[
 \boxed{\widetilde\Psi(k,t)=(\pi\alpha)^{-1/4}
 e^{-(k-k_0)^2/(2\alpha)}e^{-i\hbar k^2t/(2m)}},\qquad
 \boxed{\Pi(k,t)=\frac{e^{-(k-k_0)^2/\alpha}}{\sqrt{\pi\alpha}}}.
\]
这里 $\int\Pi(k,t)dk=1$；若改用 $p$，则概率密度为 $\Pi(p/\hbar,t)/\hbar$。
\item 作逆变换，令 $s=k-k_0$、$\tau=\alpha\hbar t/m$、$v_0=\hbar k_0/m$，
\[
 \Psi(x,t)=\frac{e^{ik_0x-i\hbar k_0^2t/(2m)}}{\sqrt{2\pi}(\pi\alpha)^{1/4}}
 \int_{-\infty}^\infty
 e^{-(1+i\tau)s^2/(2\alpha)+is(x-v_0t)}ds.
\]
使用 $\int e^{-As^2+Bs}ds=\sqrt{\pi/A}\,e^{B^2/(4A)}$（$\operatorname{Re}A>0$），
得到
\[
 \boxed{\Psi(x,t)=\left(\frac\alpha\pi\right)^{1/4}
 \frac{e^{ik_0x-i\hbar k_0^2t/(2m)}}{\sqrt{1+i\tau}}
 \exp\!\left[-\frac{\alpha(x-v_0t)^2}{2(1+i\tau)}\right]}.
\]
平方根取从 $t=0$ 连续延伸的分支。由
$|1+i\tau|=\sqrt{1+\tau^2}$、$\operatorname{Re}(1+i\tau)^{-1}=(1+\tau^2)^{-1}$，
\[
 \boxed{P(x,t)=\sqrt{\frac\alpha\pi}\frac1{\sqrt{1+\tau^2}}
       \exp\!\left[-\frac{\alpha(x-v_0t)^2}{1+\tau^2}\right]}.
\]
\item 动量概率分布不变，$\langle p\rangle=\hbar k_0$、$\Delta p=\hbar\sqrt{\alpha/2}$。
位置分布中心按 $\langle x\rangle=v_0t$ 移动，宽度
\[
 \Delta x(t)=\sqrt{\frac{1+\tau^2}{2\alpha}}
\]
随时间增大，峰高按 $(1+\tau^2)^{-1/2}$ 减小，而总概率保持为 $1$。
\end{enumerate}
}

\item 
（共30分）
一自旋为 $1 / 2$ 的粒子的归一化自旋态为 $|\lambda\rangle$ 。设自旋态 $|\lambda\rangle$ 下测量
自旋角动量 $\hat{s}_{z}$ 得到 $\frac{\hbar}{2}$ 的几率是 $9 / 25$, 测量 $\hat{s}_{x}$ 得到 $\frac{\hbar}{2}$ 的几率是 $1 / 2$ 。
\begin{enumerate}
	%\renewcommand{\labelenumi}{\arabic{enumi}.}
	% A(\Alph) a(\alph) I(\Roman) i(\roman) 1(\arabic)
	%设定全局标号series=example	%引用全局变量resume=example
	%[topsep=-0.3em,parsep=-0.3em,itemsep=-0.3em,partopsep=-0.3em]
	%可使用leftmargin调整列表环境左边的空白长度 [leftmargin=0em]
	\item
	求自旋态 $|\lambda\rangle$ 。
	\item 
	求该自旋态下自旋角动量 $\hat{s}_{y} $ 的平均值。
	
	
	
\end{enumerate}
\banswer{
%采取不同表象解得结果不一样，一般在$ \sigma_{z} $中进行。下面的答案是对于$ \sigma_{z} $而言的。
	\begin{enumerate}
		%\renewcommand{\labelenumi}{\arabic{enumi}.}
		% A(\Alph) a(\alph) I(\Roman) i(\roman) 1(\arabic)
		%设定全局标号series=example	%引用全局变量resume=example
		%[topsep=-0.3em,parsep=-0.3em,itemsep=-0.3em,partopsep=-0.3em]
		%可使用leftmargin调整列表环境左边的空白长度 [leftmargin=0em]
		\item
		在$ s_{z} $表象中求解，如果采用这种设法
		\[ 
		\ket{\lambda} =  \frac{3}{5} e^{i\phi_{1}} \commamatrix[b]{1; 0} + \frac{4}{5} e^{i\phi_{2}} \commamatrix[b]{0; 1}
		 \]
		展开系数有一个相位的不确定性，其间的关系为$ \phi_{1}=\phi_{2}+\frac{\pi}{2} + n\pi $，取值不能随意，要考虑各种情况。可以这样设
		\[ 
		|\lambda\rangle=\left(\begin{tblr}{c}
			 \dfrac{3}{5}  \\
			x+i y
		\end{tblr}\right)
		 \]
	由 $S_z$ 的测量概率及归一化，$x^2+y^2=16/25$。
	再将其投影到 $|+x\rangle=(1,1)^T/\sqrt2$，
	\[
	 \frac12=|\langle+x|\lambda\rangle|^2
	 =\frac12\left|\frac35+x+iy\right|^2
	 =\frac12\left(1+\frac65x\right).
	\]
	所以 $x=0$，进而 $y=\pm4/5$。得
	\[ 
	\left|\lambda_{1}\right\rangle=\left(\begin{tblr}{c}
		\dfrac{3}{5} \\
		\dfrac{4 i}{5}
	\end{tblr}\right)  \qquad \left|\lambda_{2}\right\rangle=
\left(\begin{tblr}{c}
	\dfrac{3}{5} \\
	-\dfrac{4 i}{5}
\end{tblr}\right)
	 \]
		
	\item 
对于归一化列矢 $(a,b)^T$，
\[
 \langle s_y\rangle=\frac\hbar2(a^*,b^*)\begin{pmatrix}0&-i\\i&0\end{pmatrix}
 \begin{pmatrix}a\\b\end{pmatrix}=\hbar\operatorname{Im}(a^*b).
\]
所以两态分别给出 $\boxed{\langle s_y\rangle_{\lambda_1}=12\hbar/25}$ 与
$\boxed{\langle s_y\rangle_{\lambda_2}=-12\hbar/25}$。
两种相对相位均满足题设，不能只保留其中一种。
		
	\end{enumerate}
	
	
}


\newpage
\item 
（共30分）两个质量均为 $m$ 的非全同粒子在一维无限深势阱
$$
V(x)=\left\{\begin{tblr}{ll}
	0, & 0\leq x \leq a \\
	\infty, & x<0\ \text{或}\ x>a
\end{tblr}\right.
$$
中运动。
\begin{enumerate}
	%\renewcommand{\labelenumi}{\arabic{enumi}.}
	% A(\Alph) a(\alph) I(\Roman) i(\roman) 1(\arabic)
	%设定全局标号series=example	%引用全局变量resume=example
	%[topsep=-0.3em,parsep=-0.3em,itemsep=-0.3em,partopsep=-0.3em]
	%可使用leftmargin调整列表环境左边的空白长度 [leftmargin=0em]
	\item
	不考虑两个粒子间的相互作用, 求该体系的能级和归一化能量本征态。 
	\item 
	设两个粒子之间的相互作用为
	$$
	H^{\prime}\left(x_{1}, x_{2}\right)=\left\{\begin{tblr}{ll}
		V, & \left(\left|x_{2}-x_{1}\right| \leq b\right) \\
		0, & \left(\left|x_{2}-x_{1}\right|>b\right)
	\end{tblr}\right.
	$$
	其中 $x_{1}$ 和 $x_{2}$ 分别为两个粒子的坐标, $b \leq a, V$ 为常数。以 $H^{\prime}$ 作为微扰, 求基态 能量的一级修正, 结果只保留至 $b / a$ 的一次项。
	
\end{enumerate}
\banswer{
\begin{enumerate}
	%\renewcommand{\labelenumi}{\arabic{enumi}.}
	% A(\Alph) a(\alph) I(\Roman) i(\roman) 1(\arabic)
	%设定全局标号series=example	%引用全局变量resume=example
	%[topsep=-0.3em,parsep=-0.3em,itemsep=-0.3em,partopsep=-0.3em]
	%可使用leftmargin调整列表环境左边的空白长度 [leftmargin=0em]
	\item
	波函数为
	\[ 
	\psi_{n_{1}, n_{2}}\left(x_{1}, x_{2}\right)=\frac{2}{a} \sin \frac{n_{1} \pi x_{1}}{a} \sin \frac{n_{2} \pi x_{2}}{a}
	 \]
	 能量为
	 \[ 
	 E_{n_{1}, n_{2}}=\frac{\pi^{2} \hbar^{2}}{2 m a^{2}}\left(n_{1}^{2}+n_{2}^{2}\right)
	  \]
	\item 
	基态能量$E_{1,1}=\frac{\pi^{2} \hbar^{2}}{m a^{2}}$，基态包函数$\psi_{1,1}^{(0)}\left(x_{1}, x_{2}\right)=\frac{2}{a} \sin \frac{\pi x_{1}}{a} \sin \frac{\pi x_{2}}{a}$，无简并，能量一级修正项为
	\[ 
	E^{(1)}=\iint \frac{2}{a} \sin \frac{\pi x_{1}}{a} \sin \frac{\pi x_{2}}{a} H^{\prime} \frac{2}{a} \sin \frac{\pi x_{1}}{a} \sin \frac{\pi x_{2}}{a} d x_{1} d x_{2}
	 \]
直接计算这个积分相当复杂，但考虑到答案只要求我们保留至$ b/a $的一次项即可，便可通过让$ b \rightarrow 0 $的形式来化简。首先化简积分区域的表达式
\[ 
E^{(1)}=V\left(\frac{2}{a}\right)^{2} \int_{0}^{a} \sin ^{2} \frac{\pi x_{1}}{a} \int_{\max(0,x_{1}-b)}^{\min(a,x_{1}+b)} \sin ^{2} \frac{\pi x_{2}}{a} d x_{2} d x_{1}
 \]
 其次，在 $b/a\ll1$ 时，只保留窄带的首项，内层积分可近似为
$2b\sin^2(\pi x_1/a)$。靠近两端的截断区对首项没有贡献，因为基态密度在端点为零。
再用 $\int_0^a\sin^4(\pi x/a)dx=3a/8$，化简为
 \[ 
\begin{aligned}
	E^{(1)}&=V\left(\frac{2}{a}\right)^{2} \int_{0}^{a} \sin ^{2} \frac{\pi x_{1}}{a} \int_{\max(0,x_{1}-b)}^{\min(a,x_{1}+b)} \sin ^{2} \frac{\pi x_{2}}{a} d x_{2} d x_{1} \\
	&\simeq V\left(\frac{2}{a}\right)^{2} \int_{0}^{a} \sin ^{2} \frac{\pi x_{1}}{a} \cdot 2b \cdot \sin ^{2} \frac{\pi x_{1}}{a} d x_{1}\\
	&=3V\frac{b}{a}
\end{aligned}
  \]
\end{enumerate}

\begin{note}
	此题，其它能级的简并很复杂。举个例子：$ 33 ^{2} +4 ^{2}=31 ^{2} +12 ^{2} = 23 ^{2}+24 ^{2}=9 ^{2}+32 ^{2} $
\end{note}
	
}

\item 
（共30分）
\begin{enumerate}
	%\renewcommand{\labelenumi}{\arabic{enumi}.}
	% A(\Alph) a(\alph) I(\Roman) i(\roman) 1(\arabic)
	%设定全局标号series=example	%引用全局变量resume=example
	%[topsep=-0.3em,parsep=-0.3em,itemsep=-0.3em,partopsep=-0.3em]
	%可使用leftmargin调整列表环境左边的空白长度 [leftmargin=0em]
	\item
	考虑由两个自旋为 $\frac{1}{2}$ 的非全同粒子组成的系统, 两个粒子的
	自旋算符分别记为 $\hat{S}_{1}$ 和 $\hat{S}_{2} $。 求系统总自旋算符 $\hat{S}_{12}^{2}$、$ \hat{S}_{12, z}\left(\hat{S}_{12}=\hat{S}_{1}+\hat{S}_{2}\right)$ 的所有共同本征态和对应的本征值。
	\item 
	考虑由三个自旋为 $\frac{1}{2}$ 的非全同粒子组成的系统，第三个粒子的自旋记为 $\hat{S}_{3}$ 。 求总自旋算符 $\hat{S}^{2}$ 、$ \hat{S}_{z}\left(\hat{S}=\hat{S}_{12}+\hat{S}_{3}\right)$ 的所有共同本征态和对应的本征值。
	
	
\end{enumerate}


\banswer{
\begin{enumerate}
	%\renewcommand{\labelenumi}{\arabic{enumi}.}
	% A(\Alph) a(\alph) I(\Roman) i(\roman) 1(\arabic)
	%设定全局标号series=example	%引用全局变量resume=example
	%[topsep=-0.3em,parsep=-0.3em,itemsep=-0.3em,partopsep=-0.3em]
	%可使用leftmargin调整列表环境左边的空白长度 [leftmargin=0em]
	\item
	共有$ 4 $个本征态，$ \left\{ S^{2}_{12},S_{12z} \right\} $对应$ \ket{lm} $态如下
	\[ 
	\begin{tblr}{ll}
	 \ket{11}&=\ket{\uparrow\uparrow}\\
	 \ket{10}&=\frac{1}{\sqrt{2}}\left( \ket{\uparrow \downarrow} + \ket{\downarrow \uparrow} \right)\\
	 \ket{1,-1}&=\ket{\downarrow\downarrow}\\
	 \ket{00}&=\frac{1}{\sqrt{2}}\left( \ket{\uparrow \downarrow} - \ket{\downarrow \uparrow} \right)
	\end{tblr}
	 \]
	\item 
	共有$ 8 $个本征态，以$ \left\{ S^{2}_{12}, S^{2},S_{z} \right\} $为完全集，对应态为$ \ket{ljm} $如下
	\[ 
	\begin{tblr}{ll}
	\ket{1, \frac{3}{2} , \frac{3}{2} }&=\ket{\uparrow\uparrow\uparrow}\\
	\ket{1, \frac{3}{2} , \frac{1}{2} }&=\frac{1}{\sqrt{3}}\left( \ket{ \downarrow\uparrow\uparrow} + \ket{\uparrow\downarrow \uparrow}+\ket{\uparrow \uparrow\downarrow} \right)\\
	\ket{1, \frac{3}{2} , -\frac{1}{2} }&=\frac{1}{\sqrt{3}}\left( \ket{ \downarrow\downarrow\uparrow} + \ket{\downarrow\uparrow\downarrow }+\ket{\uparrow \downarrow\downarrow} \right)\\
	\ket{1, \frac{3}{2} ,- \frac{3}{2} }&=\ket{\downarrow\downarrow\downarrow}\\
	\ket{0, \frac{1}{2} , \frac{1}{2} }&=\frac{1}{\sqrt{2}}\left( \ket{\uparrow \downarrow\uparrow} - \ket{\downarrow \uparrow\uparrow} \right)\\
	\ket{0, \frac{1}{2} , -\frac{1}{2} }&=\frac{1}{\sqrt{2}}\left( \ket{\uparrow \downarrow\downarrow} - \ket{\downarrow \uparrow\downarrow} \right)\\
	\ket{1, \frac{1}{2} , \frac{1}{2} }&=\frac{1}{\sqrt{6}}\left( \ket{ \downarrow\uparrow\uparrow} + \ket{\uparrow\downarrow \uparrow}-2\ket{\uparrow \uparrow\downarrow} \right)\\
	\ket{1, \frac{1}{2} , -\frac{1}{2} }&=\frac{1}{\sqrt{6}}\left( -\ket{ \downarrow\downarrow\uparrow} - \ket{\downarrow\uparrow\downarrow }+2\ket{\uparrow \downarrow\downarrow} \right)
\end{tblr}	
	 \]
两自旋情形，先由最高权态 $|\uparrow\uparrow\rangle$ 出发，用
$S_{12,-}=S_{1,-}+S_{2,-}$ 生成三重态，其正交补就是单态。
三自旋情形，先耦合 $S_{12}=0,1$，再与 $s_3=1/2$ 耦合，因此
$0\otimes\tfrac12=\tfrac12$、$1\otimes\tfrac12=\tfrac32\oplus\tfrac12$，共 $2+4+2=8$ 个态。
$S=3/2$ 四重态由 $|\uparrow\uparrow\uparrow\rangle$ 连续降阶得到；
$S_{12}=0$ 的两个态直接为前两粒子的单态乘第三粒子的上、下自旋。
剩下 $M=1/2$ 的态可设为
\[
 a|\downarrow\uparrow\uparrow\rangle
 +b|\uparrow\downarrow\uparrow\rangle
 +c|\uparrow\uparrow\downarrow\rangle.
\]
与同 $M$ 的四重态和 $S_{12}=0$ 态正交，分别给出
$a+b+c=0$、$b-a=0$；归一化又要求 $|a|^2+|b|^2+|c|^2=1$。
取 $a=b=1/\sqrt6$、$c=-2/\sqrt6$，即得表中倒数第二态；
再用总降算符得到最后一态。两粒子表中 $S_{12}^2$、$S_{12,z}$ 的本征值分别为
$S_{12}(S_{12}+1)\hbar^2$、$M_{12}\hbar$；三粒子表中 $S^2,S_z$ 的本征值分别为
$S(S+1)\hbar^2$、$M\hbar$。额外标签 $S_{12}$ 区分两个 $S=1/2$ 的副本。
\end{enumerate}
}

	
 \end{enumerate}

 
